\phantomsection\addcontentsline{toc}{chapter}{Streszczenie}
\chapter*{Streszczenie}\markboth{Streszczenie}{Streszczenie}

% \noindent
Celem niniejszej pracy jest ocena wpływu mechanizmów wirtualnych wątków i~współbieżności strukturalnej, będących częścią inicjatywy Project Loom, na wydajność oraz złożoność implementacyjną aplikacji współbieżnych w środowisku JVM. W pracy analizowane są trzy główne podejścia do współbieżności: wirtualne wątki (Project Loom), Reactive Streams oraz Kotlin Coroutines. Badanie opiera się na realistycznych scenariuszach produkcyjnych, obejmujących obsługę równoczesnych żądań HTTP, integrację z zewnętrznymi API, przetwarzanie danych w bazach danych oraz intensywne przetwarzanie strumieniowe danych telemetrycznych.

\vspace{0.3cm}
% \noindent
Praca koncentruje się zarówno na aspektach technicznych (przepustowość, opóźnienia, zużycie zasobów), jak i implementacyjnych (czytelność kodu, złożoność architektury, propagacja kontekstu i błędów). Wyniki testów pozwalają na ocenę, w~jakim stopniu wirtualne wątki upraszczają projektowanie aplikacji współbieżnych, a także na sformułowanie praktycznych rekomendacji dla zespołów developerskich rozważających migrację do Loom lub hybrydowe podejścia współbieżności w JVM.
\newpage

\chapter*{Abstract}

\noindent
\Large{\textbf{“A Comparative and Integrative Analysis of Project Loom Mechanisms and Existing Concurrency Models in the JVM Environment for Selected Application Scenarios”}}

\vspace{1cm}
\normalsize{}

\noindent
The engineering thesis describes the development of a mobile application made using the Flutter framework. The application is designed for scanning barcodes of food products to obtain detailed information about them. In addition to providing product information such as ingredients, allergens presence, and nutritional value, the application allows users to determine if a product is suitable for them based on their dietary restrictions and preferences. The application features a system for retrieving product information, which reduces the need to query external APIs. Furthermore, the application does not require an Internet connection to be functional.

\vspace{0.3cm}

\noindent
The thesis begins by addressing the need for easier access to information about food products' ingredients, especially among individuals with food intolerances. Subsequently, the thesis provides a detailed description of the items present on food products' labels, such as the structure of EAN barcodes, types of allergens and nutritional values along with standards for their presentation, and a detailed description of the Nutri-Score system. Next, it describes the technologies used in developing the application and the data structure it relies on, including information about products available in the Open Food Facts database. The thesis further focuses on the application's architecture, programming practices used during its development, including widget-based architecture with detailed descriptions of their functionality, as well as the process of retrieving and storing product information. The thesis concludes with insights into the application's development process and suggestions for future improvements.